\documentclass[11pt,oneside]{book}
\usepackage[a4paper,margin=1in]{geometry}
\usepackage{amsmath}
\usepackage{listings}
\usepackage{graphicx}
\usepackage[utf8]{inputenc}
\DeclareUnicodeCharacter{2009}{\,}
\usepackage{xcolor}
\usepackage{fvextra}   % extends fancyvrb; gives breaklines, breakanywhere, etc.
\usepackage[colorlinks=true,
            linkcolor=blue,
            urlcolor=blue,
            citecolor=blue]{hyperref}
            
\title{Inverse Adding–Doubling Tutorial}
\author{Scott Prahl}
\date{July 2025}

\usepackage{listings,xcolor}
\definecolor{codebg}{gray}{0.95}

\lstnewenvironment{grayverb}{%
  \lstset{
    basicstyle=\ttfamily\small,
    backgroundcolor=\color{codebg},
    xleftmargin=2em,
    frame=single,
    rulecolor=\color{codebg},
    framerule=0pt,
    showstringspaces=false,
    columns=fullflexible,
    keepspaces=false,
  }%
}{}

\setlength{\parindent}{0pt}
\setlength{\parskip}{6pt plus 2pt minus 1pt} % choose what you like


\newcommand\pdflatex{pdf\LaTeX}
\newcommand\iadprog{\texttt{iad}}
\newcommand\rzero{R^{\mathit{diffuse}}_\mathit{0}}
\newcommand\rstd{R^{\mathit{diffuse}}_\mathit{std}}
\newcommand\rdirect{r_s^{\mathit{direct}}}
\newcommand\rdiffuse{r_s}
\newcommand\tdirect{t_s^{\mathit{direct}}}
\newcommand\tdiffuse{t_s}
\newcommand\tcross{t_{\mathit{cross}}}
\newcommand\Psphere{P_{\mathit{sphere}}}

\begin{document}
\frontmatter
\maketitle

% create an unnumbered Abstract chapter...
\chapter*{Abstract}
\addcontentsline{toc}{chapter}{Abstract}

This tutorial guides you through installing, configuring, and running the IAD program
to extract optical properties (scattering, absorption, anisotropy) from integrating‐sphere
measurements. Examples, command‐line options, and an overview of the code structure
are provided.


\tableofcontents

\mainmatter

\chapter{Introduction}
Inverse adding-doubling (IAD) is a robust method to recover intrinsic optical parameters
from measurements of total reflection ($M_R$) and transmission ($M_T$) of a sample
measured with zero, one, or two integrating spheres. The program iteratively adjusts tissue optical properties
until simulated observables match your data.
\section{Assumptions}
\begin{itemize}
  \item The sample is a homogeneous slab between two glass slides (or nothing)
  \item Phase function restricted to Henyey–Greenstein ($g$) scattering
  \item Inputs are total reflection, total tranmission, and/or unscattered transmission.
\end{itemize}

% !TEX root = 00-tutorial.tex

\chapter{Getting Started}

This tutorial assumes you have or will obtain the \texttt{iad} executable as follows:

\section{Windows}
Download the latest precompiled Windows binary (e.g.\ \texttt{iad.exe}) from github:
\begin{grayverb}
https://github.com/scottprahl/iad
\end{grayverb}
Place the binary somewhere on your \texttt{PATH}.

\section{macOS and Linux}
Clone the source and build:
\begin{enumerate}
  \item Clone the repository:
    \begin{grayverb}
git clone https://github.com/scottprahl/iad
cd iad
    \end{grayverb}
  \item Build the program:
    \begin{grayverb}
make clean
make all
    \end{grayverb}
  \item (Optional) Install or add to your 
    \begin{grayverb}
make install
    \end{grayverb}
or move the binary into your \texttt{PATH}:
    \begin{grayverb}
cp iad /usr/local/bin
    \end{grayverb}
\end{enumerate}

\section{Testing}
Example test files (e.g., \texttt{basic4.rxt}) are provided in \texttt{test/}. Run from the iad directory where you built \texttt{iad}
\begin{grayverb}
./iad test/basic_A.rxt
\end{grayverb}
and compare the output against the reference in the \texttt{basic4-A.txt} appendix

% !TEX root = 00-tutorial.tex

\chapter{Single inversion without an \texttt{.rxt} File}

\section{No scattering, no boundaries}
When the slab has no scattering ($\mu_s=0$) the unscattered transmission matches total transmission $M_U=M_T$. This is easily inverted on the command line

\begin{grayverb}
iad -V0 -d 1.0 -t 0.325
\end{grayverb}

\begin{description}
  \item[\ttfamily -V0] silence extra output (only numeric fields printed)  
  \item[\ttfamily -d 1.0] thickness in mm  
  \item[\ttfamily -t 0.325] measured $M_T$  
\end{description}

\texttt{iad} defaults to assuming no scattering (since no light bounces back) and no integrating spheres. 

A typical one-line output:
\begin{grayverb}
1 0.0000 0.0000 0.3250 0.3250 1.1239 0.0000 0.0000 * 
\end{grayverb}

Fields are:
\begin{itemize}
  \item line number (1)
  \item $M_R$ (here 0.0000)  
  \item $M_R$ fit (here 0.0000)  
  \item $M_T$ (0.3250)  
  \item $M_T$ fit (0.3250)  
  \item recovered $\mu_a$ (1.1239 mm$^{-1}$)  
  \item reduced scattering $\mu_s'$ (0.0000 default)  
  \item anisotropy $g$ (0.0000 default)  
  \item error flag (* = success)  
\end{itemize}

\section{No scattering with boundaries}
We can add boundaries to the previous example by using \texttt{-N} and \texttt{-n}.  However, now there will be specular (unscattered) reflection from the surfaces.  To account for this, we use \texttt{-c 0} to explicitly tell the program that no unscattered reflectance is included in $M_R$.

\begin{grayverb}
iad -V0 -c 0 -N 1.5 -n 1.4 -d 1.0 -t 0.325
\end{grayverb}

\begin{description}
  \item[\ttfamily -V0] silence extra output (only numeric fields printed)  
  \item[\ttfamily -c 0] no unscattered reflectance included
  \item[\ttfamily -N 1.5] refractive index of bounding slides 
  \item[\ttfamily -n 1.4] refractive index of slab
  \item[\ttfamily -d 1.0] thickness in mm  
  \item[\ttfamily -t 0.325] measured $M_T$  
\end{description}

The defaults are no scattering (since no light bounces back) and no integrating spheres. 

A typical one-line output:
\begin{grayverb}
1 0.0000 0.0000 0.3250 0.3250 1.0401 0.0000 0.0000 r
\end{grayverb}

Fields are:
\begin{itemize}
  \item line number (1)
  \item $M_R$ (assumed 0.0000)  
  \item $M_R$ fit (0.0000)  
  \item $M_T$ (0.3250)  
  \item $M_T$ fit (0.3250)  
  \item recovered $\mu_a$ (1.0401 mm$^{-1}$)  
  \item reduced scattering $\mu_s'$ (0.0000 default)  
  \item anisotropy $g$ (0.0000 default)  
  \item error flag (* means ok)  
\end{itemize}

Here the absorption is less than for the no boundary case because the unscattered reflected light at the top and bottom boundaries is accounted for.

\section{Backscattered light from thick sample}
Since we only have one measurement $M_R$, we will need to provide some extra information.  We can extract the absorption coefficient by using the
\texttt{-F} command line option.  This time we will also remove the slides on top and bottom.

\begin{grayverb}
iad -V0 -n 1.4 -d 1.0 -F 1 -r 0.3
\end{grayverb}

\begin{description}
  \item[\ttfamily -V0] silence extra output (only numeric fields printed)  
  \item[\ttfamily -n 1.4] refractive index of slab
  \item[\ttfamily -d 1.0] thickness in mm  
  \item[\ttfamily -F 1.0] scattering coefficient in mm$^{-1}$
  \item[\ttfamily -t 0.325] measured $M_T$  
\end{description}

The defaults are no scattering (since no light bounces back) and no integrating spheres. 

A typical one-line output:
\begin{grayverb}
 1 0.3000 0.3000 0.0000 0.0000 0.0958 1.0000 0.0000 *
 \end{grayverb}

Fields are:
\begin{itemize}
  \item line number (1)
  \item $M_R$ (0.3000)  
  \item $M_R$ fit (0.3000)  
  \item $M_T$ (0.0000)  
  \item $M_T$ fit (0.0000)  
  \item recovered $\mu_a$ (0.0958 mm$^{-1}$)  
  \item reduced scattering $\mu_s'$ (specified 1.0000mm$^{-1}$)  
  \item anisotropy $g$ (0.0000 default)  
  \item error flag (* means ok)  
\end{itemize}

In this example, the scattering coefficient was kept constant and the absorption coefficient varied until the calculated value for total
reflection matched the command-line value of $M_R = 0.3$

% !TEX root = 00-tutorial.tex

\chapter{Integrating Spheres}

We ignore the unscattered transmission measurements for now.

\section{Perfect measurements of total reflection/transmission (or diffuse reflection/transmission)}

Sometimes one just wants to convert a total reflection measurement to (assuming sample thickness
is 1mm and the scattering coefficient is 1mm$^{-1}$).  Here the \texttt{-V 0} option suppresses 
extra output

\begin{grayverb}
iad -r 0.5
\end{grayverb}

to get an absorption coefficient of 0.0659mm$^{-1}$.

%If the total diffuse reflection is known then use
%\begin{grayverb}
%iad -V 0 -r 0.5 -c 0
%\end{grayverb}

\section{Dual Beam Experiments}
If you happen to have a spectrophotometer with an integrating sphere that 
alternates between illuminating the sample and a reflection standard, then 
you should use the \texttt{-X} command-line and set the number of spheres to one.  

The \texttt{-X} option assumes that your measurements have already been corrected f
or most sphere effects.  

\section{Single sphere use}
Spheres are needed to collect all light reflected or transmitted by the sample.  Most
people use a single sphere; they measure sample reflectance and then move the sample
to another port to measure the transmittance.

The normalized reflection and transmission are defined six measurements 

\begin{equation}
M_R \equiv r_\mathit{std}\cdot{{R(\rdirect, \rdiffuse) - R(0, 0)\over 
                    R(r_{\mathit{std}}, r_{\mathit{std}})- R(0, 0)} }.
\label{eqmr}
\end{equation}

and

\begin{equation}
M_T \equiv {{T(\tdirect, \rdiffuse) - T_{\mathit{dark}}\over 
                    T(0, 0) - T_{\mathit{dark}}} }
\label{eqmt}
\end{equation}

where all the quantities are defined in the figure above.  Typically one collects spectra
across all wavelengths for each these six quantities.  These are collected in a \texttt{.rxt} file as below.

\begin{figure}
\begin{center}
\includegraphics[width=0.45\textwidth]{../Ch3spheresR.pdf}
\,
\includegraphics[width=0.34\textwidth]{../Ch3spheresT.pdf}
\end{center}
\caption{Measurements needed when using a single sphere.}
\end{figure}


\begin{verbatim}
IAD1  # Must be first four characters
1.38   # Index of refraction of the sample 
1.00   # Index of refraction of the top and bottom slides 
1.62   # [mm] Thickness of sample 
0.00   # [mm] Thickness of slides 
2.00   # [mm] Diameter of illumination beam 
0.99   # Reflectivity of the reflectance calibration standard

1      # Number of spheres used for the measurement

# Properties of sphere used for reflectance measurements
101.6  # [mm] Sphere Diameter
12.7   # [mm] Sample Port Diameter
12.7   # [mm] Entrance Port Diameter
0.6    # [mm] Detector Port Diameter
0.94   # Reflectivity of the sphere wall

# Properties of sphere for transmittance measurements
101.6  # [mm] Sphere Diameter
12.7   # [mm] Sample Port Diameter
12.7   # [mm] Third Port Diameter
0.6    # [mm] Detector Port Diameter
0.94   # Reflectivity of the sphere wall

L      r         t     g
852 0.315087 0.606721 0.9
853 0.333245 0.590481 0.9

\end{verbatim}

By creating such a file, all the information about a particular experiment is captured in one place.  If this file is called
\texttt{example.rxt} then 

\begin{grayverb}
iad example.rxt
\end{grayverb}

%\input{03_rxt_files.tex}
%\input{05_scattering.tex}

\cleardoublepage
\appendix
\addcontentsline{toc}{chapter}{basic_A.rxt}
Listing of \texttt{basic_A.rxt}

\VerbatimInput[
  fontsize=\small,     % typewriter font size
  breaklines=true,     % wrap long lines
  breakanywhere=true,  % allow breaks inside long "words"
  obeytabs=true,       % respect tabs
  tabsize=4,           % ...with this width
  numbers=left,        % line numbers (optional)
  frame=single         % simple box (optional)
]{../../test/basic_A.rxt}

\cleardoublepage
\appendix
\addcontentsline{toc}{chapter}{basic_A.txt}
Listing of \texttt{basic_A.txt}

\VerbatimInput[
  fontsize=\scriptsize,     % typewriter font size
  breaklines=true,     % wrap long lines
  breakanywhere=true,  % allow breaks inside long "words"
  obeytabs=true,       % respect tabs
  tabsize=4,           % ...with this width
  numbers=left,        % line numbers (optional)
  frame=single         % simple box (optional)
]{../../test/basic_A.txt}

\end{document}
